%!TEX root = ../main.tex

\section{Related Literature}\label{sec:literature}

This paper connects to several strands of the finance literature.

\paragraph{Structural credit models.}
The structural approach to corporate liabilities originates with
\citet{merton1974pricing}, who models equity as a call option on the firm's
assets.  \citet{leland1994corporate} extends this framework to incorporate
endogenous default boundaries and optimal capital structure, while
\citet{leland1996optimal} adds debt maturity choice.  Our model shares the
structural perspective---MSTR equity is a levered claim on BTC---but
introduces an accumulation option absent from classical models, where asset
dynamics are exogenous.

\paragraph{Closed-end fund premia and discounts.}
The persistent deviation of closed-end fund (CEF) prices from NAV is a
well-documented anomaly.  \citet{lee1991investor} attribute discounts to
investor sentiment; \citet{pontiff1995closed} emphasises managerial expenses
and dividend policy; \citet{cherkes2009closed} model premia as reflecting
liquidity benefits.  The BTC treasury vehicle resembles a CEF in that equity
trades at a premium to the underlying asset value, but differs fundamentally:
the premium is sustained by an \emph{active} accumulation strategy rather
than passive holding, and the manager can issue new shares to grow the asset
base---a feature absent from standard CEFs.

\paragraph{Convertible securities and hybrid capital.}
The valuation of convertible bonds follows \citet{tsiveriotis1998valuing},
who decompose the instrument into equity and debt components.
\citet{brennan1977convertible} and \citet{ingersoll1977contingent} develop
structural models for convertible pricing.  Strategy's capital stack---common
equity, convertible notes, convertible preferred (STRK), and
non-convertible preferred (STRF, STRC, STRD, STRE)---creates a layered
priority structure that we model explicitly.  The interaction between
preferred dividends and common equity accumulation dynamics is novel.

\paragraph{Reflexivity and feedback in financial markets.}
\citet{soros1987alchemy} introduces the concept of reflexivity---the
feedback between market prices and the fundamentals they are supposed to
reflect.  In formal models, \citet{brunnermeier2014macroeconomic} and
\citet{he2013intermediary} study feedback loops between asset prices and
financial intermediary balance sheets.  \citet{brunnermeier2009market}
formalise margin spirals where declining collateral values trigger forced
sales.  Our reflexive flywheel mechanism is distinct: it operates through
\emph{voluntary} equity issuance at a premium rather than forced
deleveraging, though the death spiral shares features with margin spirals.

\paragraph{Cryptocurrency finance.}
The theoretical literature on cryptocurrency markets is growing rapidly.
\citet{biais2023equilibrium} model Bitcoin as a coordination game with
multiple equilibria in adoption.  \citet{cong2021tokenomics} develop a
dynamic model of token-based platforms.  \citet{foley2019sex} and
\citet{makarov2020trading} study Bitcoin market microstructure and
arbitrage.  On the empirical side, \citet{liu2021risks} document
cryptocurrency risk factors, and \citet{griffin2020bitcoin} examine price
manipulation.  Our contribution is orthogonal: we model the \emph{corporate
vehicle} that holds Bitcoin rather than the cryptocurrency itself, though
the market impact analysis (Section~\ref{sec:market_impact_theory}) connects
the two through the price feedback channel.

\paragraph{Bitcoin ETFs and passive vehicles.}
The approval of spot Bitcoin ETFs in the United States in January~2024
created a benchmark for pure delta-one BTC exposure.
\citet{alexander2024bitcoin} study ETF premium dynamics and the
creation/redemption mechanism.  We provide a formal comparison between the
ETF (zero accumulation option) and the treasury vehicle (positive
accumulation option), deriving conditions under which each dominates
(Section~\ref{sec:etf_comparison}).

\paragraph{Gap in the literature.}
No existing work provides a structural theory for the NAV premium of a BTC
treasury vehicle.  The CEF literature treats premia as anomalies driven by
sentiment or frictions; the structural credit literature assumes exogenous
asset dynamics; and the crypto-finance literature focuses on token design
and market microstructure.  Our paper bridges these literatures by showing
that the premium is the equilibrium price of an accumulation option whose
value depends endogenously on the premium itself.
